\documentclass{report}
\usepackage{amsmath,amssymb}
\setlength{\parindent}{0mm}
\setlength{\parskip}{1em}
\begin{document}
\begin{center}
\rule{6in}{1pt} \
{\large Robert Nourgaliev \\
{\bf On the Preconditioning of a High-Order RDG-based All-Speed Navier-Stokes Solver}}

Lawrence Livermore National Laboratory (LLNL) \\ P O Box 808 \\ L-170 \\ Livermore CA 94551
\\
{\tt nourgaliev1@llnl.gov}\\
Brian Weston\\
Robert Nourgaliev\\
Andy Anderson\\
J.-P. Delplanque\end{center}

We investigate the preconditioning of an all-speed Navier-Stokes solver,
based on the orthogonal-basis Reconstructed Discontinuous Galerkin (RDG)
space discretization, and integrated using a high-order fully-implicit
time discretization method. The work is motivated by applications in
Additive Manufacturing (AM), requiring simulations of laser-induced
powder melting with formation and subsequent solidification of liquid
metal pools, with numerous numerically challenging issues. These include
modeling of liquid-solid interfaces, due to powder
melting/solidification; gas-solid and gas-liquid multi-material
interfaces, with representation of ambient gaseous media (air, Argon),
using a fully-compressible formulation; and complex interfacial physics,
including surface tension and Marangoni effects, adequate representation
of fluid-solid-gas contacts and wetting phenomena.

Governing equations for melting/solidifying metallic substances are
discretized with (up to) the 4th-order accurate (Reconstructed) DG
method. Since we are interested in the low-speed limit, the degrees of
freedom solved for are those for pressure, velocity and
temperature/specific internal energy/enthalpy, while the non-linear
residual vectors are formulated for conservative variables (mass,
momentum and energy). The resulting set of non-linear equations is solved
using the Newton-Krylov method, with the Jacobian-free version of the
GMRES solver for linear iterations. For time discretization, we use a set
of high-order (up to the 5th) L-stable fully-implicit discretization
schemes (BDF2 and ESDIRK3-5). Due to the stiffness of the underlying
physics (combination of acoustics, thermal and viscous/material strength
effects), GMRES must be preconditioned. Our strategy for preconditioning
is a combination of a p-multigrid method with (physics-based) operator
splitting. While operator-splitting preconditioning of
(finite-volume-based) Navier-Stokes solvers has been successfully
deployed in the past, using a p-multigrid technique for the
preconditioning of DG-based solvers is relatively new, especially in the
context of the Jacobian-free Newton Krylov (JFNK) methodology. We
investigate different options of splitting governing equations for
linearized residuals, rendering a reduced set of the approximate
(preconditioning) Jacobian, which can be approximately solved for at the
preconditioning steps of the Krylov iterations, either directly or
iteratively (with Krylov- or AMG-based solvers). The figures of merit in
designing successful splitting of equations and degrees of freedom
(within p-multigrid) are � (a) the spread/clustering of eigenvalues, (b)
the condition number of the preconditioned Jacobian matrix, (c) the total
number of Krylov iterations, (d) the memory imprint due to additional
storage required for approximate Jacobian matrices, and (e) total CPU
time per non-linear iterations. The method is tested on a classical set
of benchmark problems (like lid- or thermally-driven flows), and
AM-relevant natural convection problems with embedded melting/solidifying
fronts.

This work performed under the auspices of the U.S. Department of Energy
by Lawrence Livermore National Laboratory under Contract
DE-AC52-07NA27344 and funded by the Laboratory Directed Research and
Development Program at LLNL under project tracking code 13-SI-002.


\end{document}
